
%%%%% RENAMING
%%%%%%%%%%%%%%%%%%%%%%%%%%%%%%%
%% Renaming preserves controllability 
%%%%%%%%%%%%%%%%%%%%%%%%%%%%%%%
\subsection{Renaming}
\label{subsection:renaming}
The theorem above requires the result of quotienting one of the LTS in $\mathcal{M}$ (i.e., $M_{1}^{c} \slash \! \sim_\Upsilon$) to be deterministic. This is not always the case but can be achieved by first renaming any events that may introduce non-determinism.

We use \textit{renaming} to disambiguate nondeterminism as in \cite{wong1989computation}. 

\begin{definition}[Renaming]
Let $\Sigma_1$, $\Sigma_2$ two sets of events and $\Sigma_c$ a set of controllable events. We define $\rho: \ \Sigma_2 \xrightarrow[]{} \Sigma_1$ as a \textit{renaming} which \textit{preserves controllability of labels} (i.e., if $a \in \Sigma_2$, then $\rho(a) \in \Sigma_c$  iff $a \in \Sigma_c$).
\end{definition} 

Renaming to avoid nondeterminism will introduce new events. Thus controllers with a different alphabet from that of the original plant will be synthesised. This entails that, when the plant exhibits an event $e \in \Sigma$, we have to determine which of the new events $e' \in \rho(\Sigma)$ should be triggered in the controller. This is achieved using a \textit{distinguisher}\cite{mohajerani_framework_2014}.

\begin{definition}[$\rho$-\textit{distinguisher} \cite{mohajerani_framework_2014}]
Let $\rho: \ \Sigma_2 \xrightarrow[]{} \Sigma_1$ be a renaming. An LTS $D$ with alphabet $\Sigma_2$ is a $\rho$-\textit{distinguisher} if, for all traces $s, t \in \mathcal{L}(D)$,  $\rho(s) = \rho(t)$, implies $s = t$.
\end{definition}

Renaming with $\rho$-distinguishers  guarantees that only one of the renamed events is enabled in each state. \vb{No entiendo la frase anterior. No me compila.  El renaming está hecho con una función no con un automata. para empezar, además no entiendo el significado. }


\begin{definition}[Reverse renaming]\cite{mohajerani_framework_2014}
\label{def:reverse_renaming}
Let $M = (S_{M}, \Sigma_1, \xrightarrow[]{}_M, s_{M_0})$ be an LTS, and let $\rho: \Sigma_2 \xrightarrow{} \Sigma_1$ be a renaming. Then $\rho^{-1}(M) = (S_M, \Sigma_2, \xrightarrow[]{}_{\rho^{-1}}, s_{M_0})$ where $x \xrightarrow[]{\sigma}_{\rho^{-1}} y$ iff $x \xrightarrow[]{\rho(\sigma)} y$.

\end{definition}


The following theorem states that renaming can be applied to a synthesis tuple without changing its solutions.  The proof follows from \cite{mohajerani_framework_2014}.

\begin{theorem}[Renaming preserves controllers]
    Let $\varphi$ a GR(1) formula, $\Sigma_c \subseteq \Sigma$ a set of controllable events, $\mu$ a set of uncontrollable events, $\mathcal{M} = \{ M_1, .., M_n \}$, let $D$ be a $\rho$-distinguisher such that $\rho(D)=M_1$ and $\mathcal{M'} = \{ D, \rho^{-1}(M_2), .., \rho^{-1}(M_n) \}$. Then $(\mathcal{M}, \mathcal{C}, \rho_{1}, \mu) \cong_{\Sigma_c, \varphi} (\mathcal{M'}, \rho^{-1}(\mathcal{C}) \cup \{ D \}, \rho_{1} \circ \rho, \mu)$
\end{theorem}
